% ============================================================
% CONCLUSION — Full version for NeurIPS 2026
% ============================================================

\section{Discussion and Conclusion}
\label{sec:conclusion}

We have shown that the internal weight matrices of a trained Mamba state
space model encode Granger causal structure, enabling a simple and
effective approach to causal discovery from multivariate time series. The
SSM causal score matrix $\bS = |\bW_{\text{out}} \cdot \bW_{\text{in}}|$
recovers causal graphs without any causal-specific training, group-lasso
regularization, or statistical testing---requiring only a single matrix
multiplication on the already-learned weights.

\paragraph{Theoretical contribution.}
We proved that for linear VAR($p$) processes, the SSM causal score
recovers the Granger causal structure when the model has sufficient
capacity (Theorem~\ref{thm:var_recovery}). This establishes a formal
connection between two previously disjoint fields: state space sequence
modeling and causal discovery.

\paragraph{Empirical contribution.}
On six synthetic and four real-world benchmarks spanning climate
science, finance, and neuroscience, \methodname{} achieves AUROC
0.70--1.00, outperforming both PCMCI and cLSTM across all settings.
On climate data, it outperforms Granger (0.97 vs.\ 0.93); on
financial data, it achieves perfect recovery (AUROC 1.0). Once
trained, causal extraction requires only a single matrix
multiplication ($<$1ms).

\paragraph{Limitations.}
We acknowledge four limitations:
\begin{enumerate}[leftmargin=1.5em, topsep=0pt, itemsep=0pt]
    \item \textbf{Linear theory only.} Theorem~\ref{thm:var_recovery}
    covers VAR($p$) processes. For nonlinear dynamics, our evidence is
    empirical (Lorenz system, fMRI). A nonlinear extension of the
    theory---perhaps leveraging the selectivity mechanism---is an
    important open problem.
    \item \textbf{No instantaneous causality.} SSMs are causal by
    construction ($y_t$ depends only on $x_{\leq t}$), so \methodname{}
    detects only lagged effects ($\tau \geq 1$). Instantaneous effects
    require additional structural assumptions.
    \item \textbf{Scalability to large $K$.} The $K \times K$ score
    matrix may become noisy for $K > 50$ without regularization.
    Sparse penalty on $\bW_{\text{in}}$, $\bW_{\text{out}}$ during
    training is a natural extension.
    \item \textbf{Confounding.} Like all Granger-based methods,
    \methodname{} cannot distinguish true causation from confounded
    association without additional assumptions about unobserved
    variables.
\end{enumerate}

\paragraph{Broader impact.}
\methodname{} enables faster causal analysis in domains where time series
data is abundant but causal structure is unknown---climate science,
neuroscience, economics, and genomics. By working post-hoc on pre-trained
models, it opens the door to causal analysis of existing time series
foundation models. We note that causal discovery from observational data
should be interpreted as hypothesis generation, not confirmation; any
discovered causal link should be validated through domain expertise or
intervention when possible.

\paragraph{Future work.}
Three directions are promising: (1)~extending the theory to nonlinear
SSMs by analyzing the data-dependent $\bB_t$, $\bC_t$ pathway;
(2)~developing regularized variants for high-dimensional settings
($K > 100$); and (3)~applying \methodname{} to pre-trained time series
foundation models, where the training cost is amortized across datasets
and the causal extraction becomes truly zero-cost.
